\documentclass[12pt,a4paper]{article}
\usepackage{fontspec}
\setmainfont{DejaVu Serif}
\usepackage[russian]{babel}
\usepackage{amsmath,amssymb,mathtools}
\usepackage{geometry}
\usepackage{array,booktabs,longtable}
\usepackage{xcolor}
\usepackage{tikz}
\usepackage{hyperref}
\geometry{a4paper,left=20mm,right=20mm,top=19mm,bottom=21mm}
\hypersetup{colorlinks=true,linkcolor=blue!60!black,urlcolor=blue!60!black,pdftitle={Ревью домашней работы}}
\setlength{\parindent}{0pt}
\setlength{\parskip}{6pt}
\setlength{\tabcolsep}{4pt}
\renewcommand{\arraystretch}{1.27}
\newcolumntype{L}[1]{>{\raggedright\arraybackslash}p{#1}}
\newcommand{\taskheading}[1]{\phantomsection\label{task:#1}\section*{Задание №#1}}
\newcommand{\reviewline}[1]{\par\medskip\textbf{#1}\par\smallskip}
\begin{document}

\begin{titlepage}
\thispagestyle{empty}
\begin{center}
\vspace*{0.30\textheight}
{\Huge\bfseries Ревью домашней работы\par}
\vspace{2.7cm}
{\large Автор: Лёвин Артём Александрович\par}
\vspace{0.6cm}
{\large 24 октября 2026 года\par}
\end{center}
\vfill
\begin{center}
\begin{tikzpicture}[x=1cm,y=1cm]
\draw[blue!35!black,line width=0.65pt]
  (0,0) .. controls (2.2,0.8) and (4.8,-0.8) .. (7,0)
  .. controls (8.8,0.55) and (10.3,0.32) .. (11,0);
\end{tikzpicture}
\end{center}
\end{titlepage}

\section*{Сводная таблица ошибок}
\footnotesize
\begin{longtable}{@{}L{0.12\linewidth}L{0.17\linewidth}L{0.28\linewidth}L{0.14\linewidth}L{0.17\linewidth}@{}}
\toprule
\textbf{№ задания}&\textbf{Тема}&\textbf{Суть ошибки}&\textbf{Ваш ответ}&\textbf{Правильный ответ}\\
\midrule
\endhead
\hyperref[task:2]{2}&Падение тела&В промежуточных вычислениях потерян множитель $10$; верный итог не следует из записанных действий.&$4{,}95$ м&$4{,}95$ м\\
\hyperref[task:3]{3}&Движение по окружности&В формулу подставлено $90$ сантиметров без перевода в метры; неверны число и единицы скорости.&$30$ (указаны «см»)&$3$ м/с\\
\hyperref[task:5]{5}&Высота воды; квадратное уравнение&В последнем слагаемом коэффициент $k$ взят без квадрата; ошибочно получен отрицательный дискриминант.&Не найден&$40$ с\\
\hyperref[task:6]{6}&Высота воды; отбор корня&В последнем слагаемом коэффициент $k$ взят без квадрата; время не найдено.&Не указан&$60$ с\\
\bottomrule
\end{longtable}
\normalsize

\newpage
\thispagestyle{empty}
\null\vfill
\begin{center}
{\LARGE\bfseries Разбор ошибок}
\end{center}
\vfill\null

\newpage
\taskheading{2}
\reviewline{Текст задания}
Камень падал до поверхности воды $1{,}8$ с, а после подъёма уровня воды --- $1{,}5$ с. На сколько метров изменился уровень воды? Считать $g=10$ м/с$^2$.

\reviewline{Ваше решение}
Верно записана формула пути при свободном падении из состояния покоя:
\[
 h=\frac{gt^2}{2}.
\]
Далее в записях появляются переходы
\[
\begin{aligned}
 h_1&=\frac{10\cdot 1{,}8^2}{2}
     =\frac{3{,}24}{2}=1{,}62,\\
 h_2&=\frac{10\cdot 1{,}5^2}{2}
     =\frac{2{,}25}{2}=1{,}125.
\end{aligned}
\]
Затем вычислено $1{,}62-1{,}125=0{,}495$; этот результат исправлен, и обведено число $4{,}95$.

\reviewline{Ваш ответ}
$4{,}95$ м. Конечное число верное, но приведённая цепочка вычислений содержит ошибки.

\reviewline{Ошибка}
\textcolor{red}{Ошибка:} потерян множитель $10$ при расчёте обеих высот падения.

\reviewline{Где именно возникает ошибка}
Последний правильный шаг --- подстановка $g=10$ м/с$^2$ и времени падения в формулу $h=gt^2/2$. Первый неверный переход:
\[
\underbrace{\frac{10\cdot1{,}8^2}{2}}_{\text{правильная подстановка}}
\ \not=\ \underbrace{\frac{3{,}24}{2}}_{\text{пропущено умножение на }10}.
\]
Такая же ошибка повторяется со временем $1{,}5$ с.

\reviewline{Почему это неверно}
При возведении в квадрат $1{,}8^2=3{,}24$, но это только значение квадрата времени. Числитель содержит \emph{произведение} ускорения свободного падения на квадрат времени: $10\cdot3{,}24=32{,}4$, а не $3{,}24$. Делить на два надо после умножения на десять. Аналогично $10\cdot2{,}25=22{,}5$.

Если использовать записанные неверные промежуточные значения, то получается $0{,}495$ м. Само по себе исправление этой разности на $4{,}95$ м без пересчёта обоих путей не устраняет ошибку в обосновании, даже если итог совпал с верным ответом.

\reviewline{Ключевая идея}
\textcolor{blue!60!black}{Ключевая идея:} после подъёма воды камень проходит до поверхности меньший путь. Разность путей за два времени падения и равна подъёму уровня воды. Множитель $g$ должен участвовать в каждом вычислении пути.

\reviewline{Исправление от последнего правильного шага}
Сохраняем правильно записанную формулу и исправляем арифметику:
\[
\begin{aligned}
 h_1&=\frac{10\cdot1{,}8^2}{2}=\frac{10\cdot3{,}24}{2}=16{,}2\ \text{м},\\
 h_2&=\frac{10\cdot1{,}5^2}{2}=\frac{10\cdot2{,}25}{2}=11{,}25\ \text{м}.
\end{aligned}
\]

\reviewline{Полное правильное решение}
Пусть $h_1$ --- первоначальная глубина до воды, $h_2$ --- глубина после подъёма уровня. По условию $t_1=1{,}8$ с и $t_2=1{,}5$ с. При падении из состояния покоя $h=gt^2/2$. Поэтому подъём уровня $\Delta h$ равен
\[
\begin{aligned}
\Delta h&=h_1-h_2
=\frac{g}{2}(t_1^2-t_2^2)\\
&=5(1{,}8^2-1{,}5^2)
=5(3{,}24-2{,}25)\\
&=5\cdot0{,}99=4{,}95\ \text{м}.
\end{aligned}
\]

\reviewline{Проверка}
После подъёма воды время падения уменьшилось; рассчитанная глубина действительно уменьшилась с $16{,}2$ м до $11{,}25$ м. Их разность равна $4{,}95$ м. Единица измерения --- метр.

\reviewline{Итог}
\textcolor{teal!60!black}{Итог:} уровень воды поднялся на $\boxed{4{,}95\ \text{м}}$. Конечный ответ совпал с верным, но оба промежуточных расчёта требовали исправления.

\reviewline{Задача на закрепление}
Камень, отпущенный без начальной скорости, падал до поверхности воды $2$ с. После подъёма уровня воды время падения стало $1{,}7$ с. На сколько метров поднялась вода? Принять $g=10$ м/с$^2$.

\newpage
\taskheading{3}
\reviewline{Текст задания}
Ведро с водой вращают в вертикальной плоскости на верёвке длиной $90$ см. Какова наименьшая скорость в верхней точке, при которой вода не выливается? Принять $g=10$ м/с$^2$.

\reviewline{Ваше решение}
Записано условие предельного случая, когда давление воды на дно в верхней точке обращается в ноль:
\[
\frac{v^2}{L}-g=0,\qquad v^2=gL.
\]
После указания длины $L=90$ см выполнено вычисление
\[
 v=\sqrt{10\cdot90}=\sqrt{900}=30.
\]
Около обведённого числа $30$ стоят единицы «см».

\reviewline{Ваш ответ}
Число $30$ с обозначением сантиметров; корректная единица скорости не указана.

\reviewline{Ошибка}
\textcolor{red}{Ошибка:} в произведении $gL$ объединены величины, выраженные в разных системах единиц: $g$ дано в метрах в секунду за секунду, а $L$ --- в сантиметрах.

\reviewline{Где именно возникает ошибка}
Последний правильный шаг --- равенство $v^2=gL$. Первый неверный шаг --- подстановка $L=90$ как числа без перевода сантиметров в метры:
\[
\underbrace{v^2=gL}_{\text{верная формула}},\qquad
\underbrace{v=\sqrt{10\cdot90}}_{\text{несогласованные единицы}}.
\]

\reviewline{Почему это неверно}
Ускорение $g=10$ м/с$^2$ содержит метры. Чтобы вычислить скорость в м/с по формуле $v=\sqrt{gL}$, длина $L$ тоже должна быть в метрах. Девяносто сантиметров --- это $0{,}9$ метра, а не $90$ метров. Если подставить число $90$, значение длины увеличится в сто раз, а скорость --- в десять раз.

Кроме того, скорость измеряется в метрах в секунду или сантиметрах в секунду, но не просто в сантиметрах. Сантиметр обозначает длину, а не скорость.

\reviewline{Ключевая идея}
\textcolor{blue!60!black}{Ключевая идея:} в верхней точке вода удерживается на круговой траектории, если направленное к центру ускорение не меньше ускорения свободного падения. Минимум достигается при равенстве $v^2/L=g$.

\reviewline{Исправление от последнего правильного шага}
Сохраняем условие $v^2=gL$, но предварительно переводим длину:
\[
 90\ \text{см}=0{,}9\ \text{м},\qquad
 v^2=10\cdot0{,}9=9\ \frac{\text{м}^2}{\text{с}^2}.
\]
Поскольку скорость неотрицательна,
\[
 v=\sqrt9=3\ \text{м/с}.
\]

\reviewline{Полное правильное решение}
В верхней точке центр окружности расположен внизу, поэтому центростремительное ускорение направлено вниз. На воду действуют сила тяжести $mg$ и, пока контакт сохраняется, сила давления дна, направленная к центру. Для движения по окружности необходимо
\[
 mg+N=m\frac{v^2}{L},\qquad N\ge 0.
\]
В предельном случае наименьшей скорости $N=0$. Тогда
\[
 mg=m\frac{v_{\min}^2}{L}.
\]
Сокращая положительную массу $m$, получаем $v_{\min}^2=gL$. Подставляем $L=0{,}9$ м и $g=10$ м/с$^2$:
\[
 v_{\min}=\sqrt{gL}=\sqrt{10\cdot0{,}9}
=\boxed{3\ \text{м/с}}.
\]

\reviewline{Проверка}
Проверим ускорение: $v^2/L=9/0{,}9=10$ м/с$^2$. Оно совпадает с $g$. Это и есть граничный случай; при меньшей скорости вода не смогла бы оставаться прижатой к дну в верхней точке.

\reviewline{Итог}
\textcolor{teal!60!black}{Итог:} минимальная скорость равна $\boxed{3\ \text{м/с}}$; длину верёвки обязательно нужно выразить в метрах.

\reviewline{Задача на закрепление}
Ведро с водой движется по окружности в вертикальной плоскости на верёвке длиной $1{,}6$ м. Найдите наименьшую скорость в верхней точке, при которой вода не выливается. Принять $g=10$ м/с$^2$.

\newpage
\taskheading{5}
\reviewline{Текст задания}
Высота воды в цилиндрическом баке изменяется по закону
\[
H(t)=H_0-\sqrt{2gH_0}\,kt+\frac{g}{2}k^2t^2.
\]
Через сколько секунд останется четверть первоначального объёма, если $H_0=20$ м, $g=10$ м/с$^2$, $k=1/40$?

\reviewline{Ваше решение}
Верно учтено, что четверти объёма цилиндрического бака соответствует четверть первоначальной высоты:
\[
H=\frac{H_0}{4}=\frac{20}{4}=5\ \text{м}.
\]
Затем в последнем слагаемом подстановки вместо квадрата коэффициента использован сам коэффициент:
\[
5=20-\sqrt{2\cdot10\cdot20}\,\frac{t}{40}
  +\frac{10}{2}\,\frac{t^2}{40}.
\]
Получено уравнение $t^2-4t+120=0$ и сделан вывод, что его дискриминант отрицателен. Время не найдено.

\reviewline{Ваш ответ}
Не найден: после отрицательного дискриминанта решение остановлено.

\reviewline{Ошибка}
\textcolor{red}{Ошибка:} коэффициент $k$ не возведён в квадрат в слагаемом $\frac{g}{2}k^2t^2$.

\reviewline{Где именно возникает ошибка}
Последний правильный шаг --- $H=5$ м и запись исходного закона высоты. Первый неверный шаг --- замена $k^2=(1/40)^2$ на $1/40$:
\[
\underbrace{\frac{g}{2}k^2t^2}_{\text{исходная формула}}
=5\left(\frac1{40}\right)^2t^2
=\frac{t^2}{320},
\quad\text{а не }\frac{t^2}{8}.
\]

\reviewline{Почему это неверно}
Если $k=1/40$, то $k^2=1/1600$. Квадрат относится ко всему коэффициенту: квадрат дроби равен квадрату числителя, делённому на квадрат знаменателя. Поэтому второй член с $t^2$ в записанном решении получился в сорок раз больше нужного. В результате возникло другое квадратное уравнение, не описывающее заданную высоту воды. Его отрицательный дискриминант не означает, что в исходной задаче нет подходящего времени.

\reviewline{Ключевая идея}
\textcolor{blue!60!black}{Ключевая идея:} в цилиндрическом баке объём пропорционален высоте; сначала найдём нужную высоту, затем аккуратно подставим $k^2$, решим квадратное уравнение и оставим только физически допустимое время.

\reviewline{Исправление от последнего правильного шага}
Берём уже найденную высоту $5$ м и заменяем неверное последнее слагаемое верным:
\[
\begin{aligned}
5&=20-\sqrt{400}\,\frac{t}{40}
  +5\left(\frac1{40}\right)^2t^2,\\
5&=20-\frac{t}{2}+\frac{t^2}{320}.
\end{aligned}
\]

\reviewline{Полное правильное решение}
Площадь основания цилиндрического бака постоянна, поэтому
\[
\frac{V(t)}{V_0}=\frac{H(t)}{H_0}.
\]
Если осталось $V_0/4$, то $H(t)=H_0/4=5$ м. Из исходной формулы с данными числами получаем
\[
5=20-20\cdot\frac{t}{40}
     +\frac{10}{2}\cdot\frac{t^2}{40^2}.
\]
Выполняем действия и переносим всё в одну часть:
\[
5=20-\frac{t}{2}+\frac{t^2}{320}
\quad\Longleftrightarrow\quad
0=15-\frac{t}{2}+\frac{t^2}{320}.
\]
Умножим обе части на $320$:
\[
0=4800-160t+t^2,
\qquad t^2-160t+4800=0.
\]
Разложим левую часть на множители:
\[
(t-40)(t-120)=0.
\]
Алгебраически получены два значения: $t=40$ с и $t=120$ с. Однако задача описывает процесс \emph{до опустошения} бака. Из исходного закона следует
\[
H(t)=\left(\sqrt{H_0}-\sqrt{\frac g2}\,kt\right)^2.
\]
Первое достижение нулевой высоты происходит при
\[
t_{\text{пуст}}=\frac{\sqrt{2H_0/g}}{k}
=\frac{\sqrt{2\cdot20/10}}{1/40}
=80\ \text{с}.
\]
Следовательно, физически допустимое время удовлетворяет $0\le t\le80$ с. Из двух корней подходит только $t=40$ с.

\reviewline{Проверка}
Подставим $t=40$ с в исходную формулу после правильного упрощения:
\[
H(40)=20-\frac{40}{2}+\frac{40^2}{320}
=20-20+5=5\ \text{м}.
\]
Это четверть первоначальной высоты и, следовательно, четверть первоначального объёма. Время $40$ с меньше времени опустошения $80$ с.

\reviewline{Итог}
\textcolor{teal!60!black}{Итог:} четверть первоначального объёма останется через $\boxed{40\ \text{с}}$. Корень $120$ с исключается как не соответствующий процессу до опустошения.

\reviewline{Задача на закрепление}
В цилиндрическом баке высота воды определяется законом
\[
H(t)=H_0-\sqrt{2gH_0}\,kt+\frac{g}{2}k^2t^2.
\]
Через сколько секунд останется четверть первоначального объёма, если $H_0=25$ м, $g=8$ м/с$^2$, $k=1/20$? Учитывайте физически допустимое время.

\newpage
\taskheading{6}
\reviewline{Текст задания}
Для другого цилиндрического бака действует тот же закон
\[
H(t)=H_0-\sqrt{2gH_0}\,kt+\frac{g}{2}k^2t^2.
\]
Дано: $H_0=45$ м, $g=10$ м/с$^2$, $k=1/30$. Через сколько секунд останется $1/9$ первоначального объёма? Отберите физически допустимое время.

\reviewline{Ваше решение}
Верно найдена соответствующая высота:
\[
H=\frac{H_0}{9}=\frac{45}{9}=5\ \text{м}.
\]
При подстановке, как и в предыдущей задаче, во втором квадратном слагаемом вместо $k^2$ используется $k$:
\[
5=45-\sqrt{2\cdot45\cdot10}\,\frac t{30}
 +5\cdot\frac{t^2}{30}.
\]
После преобразований появляется выражение $240-6t+t^2$. Корректное время, удовлетворяющее условию, не получено.

\reviewline{Ваш ответ}
Не указан.

\reviewline{Ошибка}
\textcolor{red}{Ошибка:} дробь $1/30$ подставлена вместо $(1/30)^2$ в последнем члене формулы высоты.

\reviewline{Где именно возникает ошибка}
Последний правильный шаг --- $H=5$ м и подстановка $\sqrt{2\cdot10\cdot45}=30$ в линейный по времени член. Первый неверный шаг --- запись $5t^2/30$ вместо
\[
\frac{10}{2}\left(\frac1{30}\right)^2t^2
=\frac{5t^2}{900}=\frac{t^2}{180}.
\]

\reviewline{Почему это неверно}
Коэффициент $k$ входит в последний член формулы во \emph{второй степени}. Если не возвести знаменатель дроби в квадрат, то коэффициент при $t^2$ увеличивается с $1/180$ до $1/6$, то есть в тридцать раз. Тогда получается другое уравнение; последующие вычисления уже не могут дать надёжный ответ на исходный вопрос.

Ещё одно важное требование содержится прямо в условии: даже после правильного решения уравнения нужно проверить, какое из найденных значений времени относится к процессу до полного опустошения бака.

\reviewline{Ключевая идея}
\textcolor{blue!60!black}{Ключевая идея:} девятая часть объёма цилиндрического бака соответствует девятой части начальной высоты. При подстановке сохраняем квадрат коэффициента $k$ и проверяем оба корня относительно времени опустошения.

\reviewline{Исправление от последнего правильного шага}
Сохраняем найденную высоту $5$ м и верно вычисленный корень $\sqrt{900}=30$:
\[
\begin{aligned}
5&=45-30\cdot\frac t{30}
  +\frac{10}{2}\left(\frac1{30}\right)^2t^2,\\
5&=45-t+\frac{t^2}{180}.
\end{aligned}
\]

\reviewline{Полное правильное решение}
Объём воды равен произведению площади основания бака на высоту. Поскольку площадь основания постоянна, условие $V(t)=V_0/9$ означает $H(t)=H_0/9=5$ м.

Подставляем параметры в закон изменения высоты:
\[
\begin{aligned}
5&=45-\sqrt{2\cdot10\cdot45}\,\frac1{30}t
       +\frac{10}{2}\left(\frac1{30}\right)^2t^2,\\
5&=45-\frac{30}{30}t+\frac{5}{900}t^2,\\
5&=45-t+\frac{t^2}{180}.
\end{aligned}
\]
Переносим всё в правую часть, затем умножаем на $180$:
\[
0=40-t+\frac{t^2}{180}
\quad\Longleftrightarrow\quad
0=7200-180t+t^2.
\]
Следовательно,
\[
t^2-180t+7200=0
\quad\Longleftrightarrow\quad
(t-60)(t-120)=0.
\]
Получили $t_1=60$ с и $t_2=120$ с. Найдём время опустошения по той же формуле, записанной в виде квадрата:
\[
H(t)=\left(\sqrt{H_0}-\sqrt{\frac g2}\,kt\right)^2.
\]
Первое достижение нулевой высоты:
\[
t_{\text{пуст}}=\frac{\sqrt{2H_0/g}}{k}
=\frac{\sqrt{2\cdot45/10}}{1/30}
=\frac3{1/30}=90\ \text{с}.
\]
Физически допустим промежуток от начала процесса до $90$ с включительно. Поэтому $t=120$ с не подходит, а $t=60$ с подходит.

\reviewline{Проверка}
\[
H(60)=45-60+\frac{60^2}{180}
=45-60+20=5\ \text{м}.
\]
Отношение $H(60)/H_0=5/45=1/9$ равно требуемой доле объёма. Кроме того, $60<90$, следовательно, бак ещё не опустел.

\reviewline{Итог}
\textcolor{teal!60!black}{Итог:} одна девятая первоначального объёма останется через $\boxed{60\ \text{с}}$. Второй алгебраический корень $120$ с физически не подходит.

\reviewline{Задача на закрепление}
Высота воды в цилиндрическом баке подчиняется закону
\[
H(t)=H_0-\sqrt{2gH_0}\,kt+\frac{g}{2}k^2t^2.
\]
Через сколько секунд останется $1/9$ первоначального объёма при $H_0=32$ м, $g=4$ м/с$^2$, $k=1/24$? Выберите физически допустимое время.

\vfill
\begin{center}
\footnotesize Лёвин Артём Александрович эксклюзивно для Ксении Клыковой
\end{center}
\end{document}